\section{Exploratory behavioral analysis}
\label{sec:behavior}
\paragraph{Sampling and coding.}
After the experiments, we froze a codebook and sampling rule before annotation. One run was sampled from every family and condition cell, giving 52 runs, with eight or nine per configuration. The rule used design labels and random repetition selection, without outcomes or trace availability; no run was replaced. All selected traces and implementations were available. Two isolated LLM sessions read every packet separately, with final scores and model labels withheld. Both used the recorded OpenAI alias \texttt{gpt-6-astra} with \texttt{xhigh} reasoning (Appendix~\ref{app:behavior}). Packets contained the task, a source procedure for comparison, initial and final implementations, the patch, and chronological visible tool interactions. Interfaces and public feedback remained visible.

The dimensions distinguish explicit assumption recognition in the visible trace \emph{before editing} (A), retention of a recognizable source procedure without its target check (P), a concrete target protection (T), and stopping without a meaningful implementation (Z). Each decision includes event or implementation line anchors. Labels are yes, no, uncertain, and unavailable. The passes could not inspect each other's output; disagreements were retained without outcome based adjudication.

\begin{table}[t]
\centering\small\begin{tabular}{p{6.1cm}rrrr}
\toprule
Observable dimension & Pass A: Y & Pass B: Y & Both Y & Agreement\\
\midrule
A: Explicit assumption recognition before editing & 8 & 7 & 7 & 49/52\\
P: Source procedure without the target check & 20 & 22 & 20 & 50/52\\
T: Target protection present & 22 & 23 & 22 & 49/52\\
Z: No meaningful implementation & 7 & 7 & 7 & 52/52\\
\bottomrule
\end{tabular}

\caption{Two model assisted coding passes over the same 52 sampled runs. ``Both Y'' requires two positive labels; agreement includes uncertain and unavailable labels. Counts describe the sample, not behavioral prevalence across all 624 runs. Full label distributions and all eight disagreements are preserved in Appendix~\ref{app:behavior}.}
\label{tab:coding}
\end{table}

\paragraph{Visible statements and implemented protection are distinct observations.}
Both passes identify explicit assumption recognition before editing in \CodeABothYes{} runs, procedure retention without the check in \CodePBothYes{}, concrete protection in \CodeTBothYes{}, and no meaningful implementation in \CodeZBothYes{} (Table~\ref{tab:coding}). All seven runs with agreed explicit assumption recognition also have agreed protection. Fifteen further runs have agreed protection without that agreed visible statement. The visible trace therefore does not provide a complete account of when protections enter the implementation. These codes do not identify whether memory influenced an edit, particularly in conditions where no source memory was delivered.

Luna F08/B/r2 (P027) illustrates the timing distinction: its patch escapes the label, but its explanation follows the edit; both checks pass. Devstral X20/C/r1 (P015) states the changed assumption before editing, distinguishing repeatable tags from a control field that remains singular. It rejects duplicate controls and uses one interpretation for authorization and dispatch; both checks pass. Terra X28/B/r1 (P043) states that a positive prefilter result is insufficient and adds an authoritative check. Functionality passes, but an error path interrupts focal witness measurement, leaving that outcome unknown.

\paragraph{Matched examples of unchecked reuse and adaptation.}
Appendix~\ref{app:behavior} gives the matched transcript, patch, and outcome panels. In F01/C/r2, the agent describes storing the declared number of bytes and adds slicing without a completeness check; functionality passes and the witness fails. The other C repetition adds that check and passes both. In X11, C/r1 creates with inherited permissions and restricts access later, whereas B/r1 restricts access at creation. Both are functional; only B passes the creation observation. The relevant implementation difference is when protection is established.

In F08, C/r1 explicitly notes that the target value is no longer guaranteed numeric and escapes the label. N/r1 extends direct interpolation. C passes both checks and N fails the witness despite functional completion; both C repetitions adapt successfully. The agent's stated rationale connects the source assumption to the target repair. These selected comparisons explain implementation differences, without estimating a family level memory effect.

\paragraph{One failure can have different implementation mechanisms.}
All 24 Codex F02 runs complete functionally and fail the focal witness. Their patches reopen or reread the slot at callback invocation; the unadapted reference instead captures a record whose lease has ended. The same endpoint, reached by different implementations. Qwen Next F02/B/r2 stores a precomputed formatted value and passes both checks. Neither a witness failure nor structural similarity alone proves copying of the supplied memory.
